\documentclass[12pt]{article}
\usepackage[utf8]{inputenc}
\usepackage[T1]{fontenc}
\usepackage[margin=1in]{geometry}
\usepackage{setspace}
\usepackage{xcolor}
\usepackage[colorlinks=true, urlcolor=blue!60!black, linkcolor=blue!60!black]{hyperref}
\setlength{\parindent}{2em}
\setlength{\parskip}{1.25em}
\renewcommand{\baselinestretch}{1.0}
\begin{document}
\begin{center}
{\LARGE \textit{Tandon v. Newsom}, 593 U.S. 61 (2021)\par}
\vspace{0.8em}
{\large Case Brief\par}
\end{center}
\vspace{1em}

\section{Memory Jogger}
In a 5--4 emergency-docket \textit{per curiam} order, the Supreme Court enjoined California's three-household cap on in-home religious gatherings, formally adopting the ``most-favored-nation'' rule that treating \textit{any} comparable secular activity more favorably than religious exercise triggers strict scrutiny under the Free Exercise Clause.

\section{Detailed Case Facts}
During the COVID-19 pandemic, California Governor Gavin Newsom and state public-health officials enacted a tiered ``Blueprint for a Safer Economy'' restricting public and private gatherings to curb viral transmission (\href{https://www.supremecourt.gov/opinions/20pdf/20a151_4g15.pdf}{Supreme Court Slip Opinion}; \href{https://www.courtlistener.com/opinion/4872486/tandon-v-newsom/}{CourtListener}). Under California's private-gathering rules, both religious and secular gatherings inside private homes were capped at no more than three households at a time (\href{https://www.supremecourt.gov/opinions/20pdf/20a151_4g15.pdf}{Slip Opinion at 3, 5}). Simultaneously, California's Blueprint system permitted numerous secular commercial activities outside the home---including hair salons, retail stores, personal care services, movie theaters, private suites at sporting events and concerts, and indoor restaurants---to bring together members of more than three households at once under sector-specific safety protocols (\href{https://www.supremecourt.gov/opinions/20pdf/20a151_4g15.pdf}{Slip Opinion at 3}).

On October 13, 2020, a coalition of California residents---including business owners, political candidates such as Ritesh Tandon, and religious plaintiffs Pastor Jeremy Wong and Karen Busch---sued state officials challenging the COVID-19 restrictions (\href{https://caselaw.findlaw.com/court/us-dis-crt-n-d-cal-san-jos-div/2110794.html}{District Court Opinion}; \href{https://en.wikipedia.org/wiki/Tandon_v._Newsom}{Case Overview}). Prior to the pandemic, Wong and Busch regularly hosted Bible studies, theological discussions, collective prayer, and musical prayer in their homes for eight to twelve people at a time (\href{https://en.wikipedia.org/wiki/Tandon_v._Newsom}{Case Overview}). They alleged that California's three-household cap violated the Free Exercise Clause of the First Amendment (incorporated against the States via the Fourteenth Amendment) by restricting in-home worship more severely than larger gatherings in secular commercial venues (\href{https://law.justia.com/cases/federal/appellate-courts/ca9/21-15228/21-15228-2021-03-30.html}{Ninth Circuit Opinion}; \href{https://www.supremecourt.gov/opinions/20pdf/20a151_4g15.pdf}{Slip Opinion at 3}). Shortly after the plaintiffs filed their emergency application in the Supreme Court, California officials announced changes to the challenged policy scheduled to take effect on April 15, 2021, while leaving the existing restrictions in force until then and retaining authority to reimpose stricter rules (\href{https://www.supremecourt.gov/opinions/20pdf/20a151_4g15.pdf}{Slip Opinion at 3}).

\section{Procedural History}
\begin{enumerate}
\item \textbf{U.S. District Court for the Northern District of California:} After a video hearing on December 18, 2020, Judge Lucy H. Koh denied the plaintiffs' motion for a preliminary injunction on February 5, 2021 (\textit{Tandon v. Newsom}, 517 F. Supp. 3d 922 (N.D. Cal. 2021)) (\href{https://caselaw.findlaw.com/court/us-dis-crt-n-d-cal-san-jos-div/2110794.html}{FindLaw}; \href{https://en.wikipedia.org/wiki/Tandon_v._Newsom}{Case Overview}). Relying on uncontested testimony from California's public-health experts, the district court found that private in-home gatherings posed greater viral risks than commercial establishments and concluded that the three-household limit applied equally to all in-home gatherings, religious and secular alike (\href{https://www.supremecourt.gov/opinions/20pdf/20a151_4g15.pdf}{Slip Opinion, Kagan, J., dissenting, at 2}; \href{https://lawreview.gmu.edu/print__issues/religious-and-secular-comparators/}{George Mason Law Review}).
\item \textbf{U.S. Court of Appeals for the Ninth Circuit:} Plaintiffs appealed and sought an emergency injunction pending appeal. On March 30, 2021, a divided Ninth Circuit panel (Judges Milan D. Smith, Jr., and Bridget S. Bade in the majority; Judge Patrick J. Bumatay dissenting in part) denied the injunction by a 2--1 vote (\textit{Tandon v. Newsom}, 992 F.3d 916 (9th Cir. 2021)) (\href{https://law.justia.com/cases/federal/appellate-courts/ca9/21-15228/21-15228-2021-03-30.html}{Justia Ninth Circuit}; \href{https://en.wikipedia.org/wiki/Tandon_v._Newsom}{Case Overview}). The panel held that in-home religious gatherings were properly compared to in-home secular gatherings rather than to commercial venues in public buildings, applied rational-basis review, and noted that commercial precautions might not ``translate readily'' to private homes (\href{https://www.supremecourt.gov/opinions/20pdf/20a151_4g15.pdf}{Slip Opinion at 3}; \href{https://lawreview.gmu.edu/print__issues/religious-and-secular-comparators/}{George Mason Law Review}).
\item \textbf{Supreme Court of the United States:} The religious plaintiffs submitted an emergency application for injunctive relief (No. 20A151) to Justice Kagan, who referred it to the full Court (\href{https://www.supremecourt.gov/opinions/20pdf/20a151_4g15.pdf}{Slip Opinion at 1}). On April 9, 2021, without oral argument or full merits briefing, the Court granted an injunction pending appeal in a 5--4 \textit{per curiam} decision (\href{https://www.courtlistener.com/opinion/4872486/tandon-v-newsom/}{CourtListener}; \href{https://www.supremecourt.gov/opinions/20pdf/20a151_4g15.pdf}{Slip Opinion at 1, 4}).
\item \textbf{Subsequent Disposition:} On June 24, 2021, by agreement of the parties, the district court entered a stipulated final judgment granting a permanent injunction and awarding \$500,000 in reasonable attorney's fees to the plaintiffs (\href{https://en.wikipedia.org/wiki/Tandon_v._Newsom}{Case Overview}).
\end{enumerate}

\section{Judicial Votes}
\begin{itemize}
\item \textbf{Vote Breakdown:} 5--4 granting the application for injunctive relief pending appeal (\href{https://www.supremecourt.gov/opinions/20pdf/20a151_4g15.pdf}{Slip Opinion at 1, 4--5}; \href{https://www.scotusblog.com/2021/04/tandon-steals-fultons-thunder-the-most-important-free-exercise-decision-since-1990/}{SCOTUSblog}).
\item \textbf{Per Curiam Majority (5):} Justices Thomas, Alito, Gorsuch, Kavanaugh, and Barrett (unsigned \textit{per curiam} opinion; no separate concurrences) (\href{https://www.scotusblog.com/2021/04/tandon-steals-fultons-thunder-the-most-important-free-exercise-decision-since-1990/}{SCOTUSblog}).
\item \textbf{Unauthored Dissent (1):} Chief Justice Roberts noted without opinion that he would deny the application (\href{https://www.supremecourt.gov/opinions/20pdf/20a151_4g15.pdf}{Slip Opinion at 4}).
\item \textbf{Dissenting Opinion (3):} Justice Kagan filed a dissenting opinion, joined by Justices Breyer and Sotomayor (\href{https://www.supremecourt.gov/opinions/20pdf/20a151_4g15.pdf}{Slip Opinion, Kagan, J., dissenting, at 1}).
\end{itemize}

\section{Holding}
\begin{itemize}
\item \textbf{Judgment:} Application for injunctive relief pending appeal granted; California was enjoined from enforcing its three-household cap on in-home gatherings against the applicants' in-home religious exercise pending disposition of the appeal in the Ninth Circuit and any timely certiorari petition (\href{https://www.supremecourt.gov/opinions/20pdf/20a151_4g15.pdf}{Slip Opinion at 1, 4}).
\item \textbf{Rule Announced:} Four principles govern Free Exercise challenges under \textit{Employment Division v. Smith}, 494 U.S. 872 (1990): (1) government regulations are not neutral and generally applicable, and therefore trigger strict scrutiny, whenever they treat \textit{any} comparable secular activity more favorably than religious exercise---even if other comparable secular activities are treated equally or worse; (2) comparability is judged solely against the asserted government interest (the risks posed), not the reasons why people gather; (3) under strict scrutiny, the government bears the burden of showing that less restrictive precautions permitted for comparable secular activities would not suffice for religious exercise; and (4) modifying or withdrawing a challenged restriction during litigation does not moot the case or bar emergency relief where applicants remain under a constant threat of reinstatement (\href{https://www.supremecourt.gov/opinions/20pdf/20a151_4g15.pdf}{Slip Opinion at 1--3}).
\item \textbf{Holding vs.\ Dicta:} The holding encompasses the four-part comparator and strict-scrutiny framework, its application to California's three-household cap relative to less-restricted commercial venues, and the voluntary-cessation mootness ruling, because each element was necessary to find likelihood of success on the merits and grant emergency injunctive relief (\href{https://www.supremecourt.gov/opinions/20pdf/20a151_4g15.pdf}{Slip Opinion at 1--4}). The \textit{per curiam}'s concluding admonition that this was the fifth time the Court had summarily rejected the Ninth Circuit's COVID-19 Free Exercise analysis, along with its broad characterization of California's entire Blueprint System as riddled with exceptions, functions as institutional critique and dicta rather than an independent legal holding (\href{https://www.supremecourt.gov/opinions/20pdf/20a151_4g15.pdf}{Slip Opinion at 4}).
\end{itemize}

\section{Analysis of Opinions}

\subsection{Per Curiam Opinion (Justices Thomas, Alito, Gorsuch, Kavanaugh, and Barrett)}
The unsigned \textit{per curiam} opinion consolidated the Court's earlier COVID-19 religious-liberty rulings---including \textit{Roman Catholic Diocese of Brooklyn v. Cuomo}, 592 U.S. 14 (2020), and \textit{South Bay United Pentecostal Church v. Newsom}, 592 U.S. 918 (2021)---into a four-point Free Exercise framework (\href{https://www.supremecourt.gov/opinions/20pdf/20a151_4g15.pdf}{Slip Opinion at 1--3}).

First, the Court held that regulations lose neutrality and general applicability---triggering strict scrutiny---whenever they treat \textit{any} comparable secular activity more favorably than religious exercise. Treating secular in-home gatherings identically or worse than religious gatherings was ``no answer'' (\href{https://www.supremecourt.gov/opinions/20pdf/20a151_4g15.pdf}{Slip Opinion at 1}). Second, comparability turns solely on the risks various activities pose to the asserted government interest (curbing viral spread), not on the social or economic reasons why people gather (\href{https://www.supremecourt.gov/opinions/20pdf/20a151_4g15.pdf}{Slip Opinion at 2}). Third, under strict scrutiny, narrow tailoring requires the State to demonstrate that religious exercise remains more dangerous than permitted secular activities even when identical precautions are applied (\href{https://www.supremecourt.gov/opinions/20pdf/20a151_4g15.pdf}{Slip Opinion at 2}). Fourth, California's announced policy revision for April 15 did not moot the dispute because the prior cap remained in force in the interim and officials with a track record of changing rules retained authority to reinstate heightened restrictions (\href{https://www.supremecourt.gov/opinions/20pdf/20a151_4g15.pdf}{Slip Opinion at 2--3}).

Applying these standards, the majority faulted the Ninth Circuit for rejecting commercial comparators (hair salons, retail stores, personal care services, movie theaters, private suites, and indoor restaurants) merely because prior decisions involved public rather than private buildings, and for speculating that commercial precautions might not translate readily to homes (\href{https://www.supremecourt.gov/opinions/20pdf/20a151_4g15.pdf}{Slip Opinion at 3}). Because applicants were likely to succeed on the merits, faced irreparable harm from even brief losses of Free Exercise rights, and the State failed to show public health would be imperiled by less restrictive precautions, the Court granted the injunction (\href{https://www.supremecourt.gov/opinions/20pdf/20a151_4g15.pdf}{Slip Opinion at 3--4}).

\subsection{Concurring Opinions}
No separate concurring opinions were filed (\href{https://www.supremecourt.gov/opinions/20pdf/20a151_4g15.pdf}{Slip Opinion at 1--4}).

\subsection{Unauthored Dissent of Chief Justice Roberts}
Chief Justice Roberts dissented without opinion and did not join Justice Kagan's dissent, noting only that he would deny the application (\href{https://www.supremecourt.gov/opinions/20pdf/20a151_4g15.pdf}{Slip Opinion at 4}). Consistent with his prior concurrences in \textit{South Bay United Pentecostal Church v. Newsom}, Chief Justice Roberts had favored substantial deference to politically accountable state officials navigating public-health emergencies unless a state singled out houses of worship compared to similar public venues (\href{https://lawreview.gmu.edu/print__issues/religious-and-secular-comparators/}{George Mason Law Review}).

\subsection{Dissenting Opinion of Justice Kagan (Joined by Justices Breyer and Sotomayor)}
Justice Kagan argued that the Free Exercise Clause requires a State to treat religious conduct as well as comparable secular conduct, and that the obvious comparator for at-home religious gatherings was at-home secular gatherings---both of which California capped identically at three households (\href{https://www.supremecourt.gov/opinions/20pdf/20a151_4g15.pdf}{Slip Opinion, Kagan, J., dissenting, at 1}). Demanding that California equate private home gatherings with hardware stores and hair salons, she contended, forced the State to treat unlike activities equivalently (\href{https://www.supremecourt.gov/opinions/20pdf/20a151_4g15.pdf}{Slip Opinion, Kagan, J., dissenting, at 1--2}).

Even under the majority's broader comparative frame, Justice Kagan argued that the \textit{per curiam} defied the factual record (\href{https://www.supremecourt.gov/opinions/20pdf/20a151_4g15.pdf}{Slip Opinion, Kagan, J., dissenting, at 1}). Contrary to the majority's claim that the Ninth Circuit never found commercial venues less risky, Judges Smith and Bade---relying on undisputed district court findings drawn from California's public-health experts---identified three concrete risk distinctions: (1) social gatherings in homes involve longer interactions and prolonged conversations; (2) private houses are typically smaller and less ventilated than commercial buildings; and (3) masking and distancing are less likely and harder to enforce in private homes (\href{https://www.supremecourt.gov/opinions/20pdf/20a151_4g15.pdf}{Slip Opinion, Kagan, J., dissenting, at 2}). Disregarding those undisputed findings, she concluded, improperly impaired the State's scientific response to a public-health emergency (\href{https://www.supremecourt.gov/opinions/20pdf/20a151_4g15.pdf}{Slip Opinion, Kagan, J., dissenting, at 2}).

\section{Examples: Future Applications}
\begin{enumerate}
\item \textbf{Same-Side (Strict Scrutiny Triggered --- Religious Exemption Likely Required):} A city police department prohibits officers from wearing beards to promote uniform appearance, but grants medical exemptions for officers with a skin condition (pseudofolliculitis barbae) while denying exemptions to Muslim officers whose faith requires a beard. Under \textit{Tandon}, the medical exception undermines the department's asserted interest in uniformity in the exact same way as a religious beard, triggering strict scrutiny (\href{https://lawreview.gmu.edu/print__issues/religious-and-secular-comparators/}{George Mason Law Review}).
\item \textbf{Same-Side (Strict Scrutiny Triggered --- Religious Exemption Likely Required):} A public university enforces an ``all-comers'' nondiscrimination rule to strip recognition from a religious student club that requires leaders to affirm its faith tenets on sexuality, while permitting secular affinity clubs (such as women's leadership groups) to restrict membership by protected traits. Under \textit{Tandon}, exempting comparable secular student groups triggers strict scrutiny against denying an exemption to the religious club (\href{https://lawreview.gmu.edu/print__issues/religious-and-secular-comparators/}{George Mason Law Review}).
\item \textbf{Opposite-Side (Neutral and Generally Applicable --- No Religious Exemption Required):} A municipality enforces a universal residential noise ordinance barring amplified sound above 65 decibels between 10:00 p.m.\ and 7:00 a.m.\ with zero secular exemptions for concerts, parties, or rallies. A religious group seeking to broadcast amplified midnight prayers cannot trigger strict scrutiny under \textit{Tandon} because no comparable secular activity receives more favorable treatment (\href{https://www.supremecourt.gov/opinions/20pdf/20a151_4g15.pdf}{Slip Opinion at 1--2}).
\item \textbf{Opposite-Side (Non-Comparable Secular Activity --- Strict Scrutiny Not Triggered):} A state prohibits private possession of venomous snakes for public safety, while exempting accredited zoos staffed by licensed specialists under inspected containment protocols. A religious sect seeking to handle venomous snakes in private worship fails \textit{Tandon}'s comparability threshold if the record shows accredited zoos pose a materially lower safety risk than unmonitored private handling (\href{https://lawreview.gmu.edu/print__issues/religious-and-secular-comparators/}{George Mason Law Review}).
\item \textbf{Fence-Sitter (Close Call):} A state requires all hospital healthcare workers to receive a COVID-19 vaccine, granting a narrow exception for physician-certified medical contraindications but no religious exemption. Lower courts have divided sharply under \textit{Tandon}: one view holds that an unvaccinated worker poses the identical transmission risk whether exempted for medical or religious reasons (triggering strict scrutiny), while the competing view holds that medical exemptions are objectively verifiable, far less numerous, and directly advance the overarching interest in worker health, making them non-comparable (\href{https://lawreview.gmu.edu/print__issues/religious-and-secular-comparators/}{George Mason Law Review}).
\end{enumerate}

\section{Critique}
\textit{Tandon v. Newsom} has been called ``the most important free exercise decision since 1990'' because it formally adopted the ``most-favored-nation'' (MFN) interpretation of general applicability---long advocated by Professor Douglas Laycock---without formally overruling \textit{Employment Division v. Smith} (\href{https://www.scotusblog.com/2021/04/tandon-steals-fultons-thunder-the-most-important-free-exercise-decision-since-1990/}{SCOTUSblog}; \href{https://lawreview.gmu.edu/print__issues/religious-and-secular-comparators/}{George Mason Law Review}).

\textbf{Originalist and Textualist Perspectives:} Proponents of robust Free Exercise protection praise \textit{Tandon} for preventing governments from implicitly subordinating constitutionally protected religious exercise to favored secular commerce (\href{https://lawreview.gmu.edu/print__issues/religious-and-secular-comparators/}{George Mason Law Review}). Yet methodological and originalist observers note that \textit{Tandon} reshaped constitutional doctrine in a four-page emergency-docket order containing no historical analysis of Founding-era Free Exercise meaning (\href{https://www.scotusblog.com/2021/04/tandon-steals-fultons-thunder-the-most-important-free-exercise-decision-since-1990/}{SCOTUSblog}; \href{https://scholarship.law.tamu.edu/cgi/viewcontent.cgi?article=1410&context=lawreview}{Texas A\&M Law Review}). Moreover, textualist and historical scholars remain divided over whether the First Amendment's text (``make no law\ldots prohibiting the free exercise'') originally entailed conduct exemptions or barred only targeted prohibitions (\href{https://lawreview.gmu.edu/print__issues/religious-and-secular-comparators/}{George Mason Law Review}).

\textbf{Progressive and Doctrinal Perspectives:} Critics such as Jim Oleske, Alan Brownstein, Vikram Amar, and Nelson Tebbe contend that the MFN approach hollows out \textit{Smith} and conflicts with \textit{City of Boerne v. Flores}, 521 U.S. 507 (1997), because nearly every modern statute contains secular exceptions, potentially subjecting ordinary public-health, labor, and antidiscrimination laws to near-fatal strict scrutiny (\href{https://www.scotusblog.com/2021/04/tandon-steals-fultons-thunder-the-most-important-free-exercise-decision-since-1990/}{SCOTUSblog}; \href{https://verdict.justia.com/2021/04/30/exploring-the-meaning-of-and-problems-with-the-supreme-courts-apparent-adoption-of-a-most-favored-nation-approach-to-protecting-religious-liberty-under-the-free-exercise-c}{Justia Verdict}; \href{https://harvardlawreview.org/print/vol-137/a-law-unto-himselfemp-div-v-smith-494-u-s-872-879-1990-quoting-reynolds-v-united-states-98-u-s-145-167-1879-free-exercise-unequal-value-and-the-future-of-public/}{Harvard Law Review}). Doctrinal scholars further warn that \textit{Tandon} leaves unanswered how to frame the level of generality of state interests, allocate the evidentiary burden on comparability, or weigh aggregate numerosity, inviting polarized lower-court rulings (\href{https://lawreview.gmu.edu/print__issues/religious-and-secular-comparators/}{George Mason Law Review}). Procedural critics such as Steve Vladeck also fault the Court for issuing a landmark doctrinal shift on the emergency ``shadow docket'' without merits briefing, oral argument, or deference to district court fact-finding (\href{https://en.wikipedia.org/wiki/Tandon_v._Newsom}{Case Overview}; \href{https://lawreview.gmu.edu/print__issues/religious-and-secular-comparators/}{George Mason Law Review}).

\textbf{Underlying Competing Values:} \textit{Tandon} highlights the clash between equal religious liberty and democratic public-health governance. On one side is the constitutional commitment that religious conscience may not be quietly devalued below secular interests during crises; on the other are institutional competence, judicial restraint, and the practical capacity of states to design risk-calibrated regulations without triggering strict scrutiny whenever a law accommodates a single secular necessity (\href{https://lawreview.gmu.edu/print__issues/religious-and-secular-comparators/}{George Mason Law Review}; \href{https://firstamendment.mtsu.edu/article/tandon-v-newsom/}{First Amendment Encyclopedia}).

\section{Key Quotations}
\begin{quote}
``First, government regulations are not neutral and generally applicable, and therefore trigger strict scrutiny under the Free Exercise Clause, whenever they treat \textit{any} comparable secular activity more favorably than religious exercise.''
\end{quote}
\noindent--- \textit{Per Curiam}, 593 U.S. at 62 (\href{https://www.supremecourt.gov/opinions/20pdf/20a151_4g15.pdf}{Supreme Court Slip Opinion at 1}; \href{https://www.courtlistener.com/opinion/4872486/tandon-v-newsom/}{CourtListener}).

\begin{quote}
``Comparability is concerned with the risks various activities pose, not the reasons why people gather.''
\end{quote}
\noindent--- \textit{Per Curiam}, 593 U.S. at 62 (\href{https://www.supremecourt.gov/opinions/20pdf/20a151_4g15.pdf}{Supreme Court Slip Opinion at 2}).

\begin{quote}
``The State cannot `assume the worst when people go to worship but assume the best when people go to work.' ''
\end{quote}
\noindent--- \textit{Per Curiam}, 593 U.S. at 63 (quoting \textit{Roberts v. Neace}, 958 F.3d 409, 414 (6th Cir. 2020)) (\href{https://www.supremecourt.gov/opinions/20pdf/20a151_4g15.pdf}{Supreme Court Slip Opinion at 3}).

\begin{quote}
``As the per curiam's reliance on separate opinions and unreasoned orders signals, the law does not require that the State equally treat apples and watermelons.''
\end{quote}
\noindent--- Justice Kagan, dissenting, 593 U.S. at 65 (\href{https://www.supremecourt.gov/opinions/20pdf/20a151_4g15.pdf}{Supreme Court Slip Opinion, Kagan, J., dissenting, at 1}).

\end{document}
